\clearpage\begingroup\SourceLeaf{362b}\ArithmeticHead{31}{R}{DIE GRUNDLAGEN DER GEOMETRIE.}
\renewcommand{\thefootnote}{\arabic{footnote})}\setcounter{footnote}{0}
\begin{center}
\textls[100]{B. Der Ausbau der nichteuklidischen Geometrie (seit 1817).}\par\medskip
\textbf{7.}\par\smallskip
\textls[100]{Die Zeit der Geodäsie und die Flächentheorie;}\\
\textls[100]{\textsc{Schweikart} und \textsc{Taurinus} (1817—1831).}
\end{center}
\noindent\hspace*{\parindent}Die Andeutungen über die Unbeweisbarkeit des Parallelenaxioms in der
Anzeige vom Jahre 1816 hatten nicht den von \textsc{Gauss} erwarteten Erfolg gehabt,
und er hatte, das Geschrei der Böoter scheuend, sich entschlossen, bei
Lebzeiten nichts über seine Ansichten bekannt zu machen. Um so mehr
sehen wir ihn überrascht und erfreut, wenn er auf seinem einsamen Wege
Gleichstrebende antrifft. Das ereignete sich mit \textsc{Schweikart} und \textsc{Taurinus}.
\par
Der Rechtsgelehrte \textsc{Schweikart} (1780—1859) hatte 1807 eine Schrift zur
Parallelentheorie veröffentlicht\footnote{F. C. \textsc{Schweikart}, \textit{Die Theorie der Parallellinien, nebst dem Vorschlage ihrer Verbannung aus der Geometrie}, Jena und Leipzig 1807. Vgl. P. Th. S. 243—246.}, in der er beanstandete, dass man bei der
üblichen Erklärung der Parallelen als einander nicht schneidender Geraden
das Unendliche hereinziehe, und forderte, man solle beim Aufbau der Geometrie
von der Existenz der Quadrate ausgehen\footnote{In ähnlicher Weise war \textsc{Clairaut}, \textit{Éléments de Géométrie}, Paris 1741, davon ausgegangen, dass das Vorhandensein von Rechtecken durch die Anschauung gegeben sei, und hatte daraus mit grosser Klarheit die Sätze des ersten Buches der \textsc{Euklid}ischen Element abgeleitet.}. Später, zwischen 1812 und 1816,
hatte er »ohne Hilfe des elften \textsc{Euklid}ischen Axioms eine Geometrie, die er
\textls[100]{Astralgeometrie} nannte, entwickelt« (Brief von \textsc{Gerling} an W. \textsc{Bolyai} vom
31. Oktober 1854, P. Th. S. 243) und, nachdem er 1816 aus Charkow an die
\clearpage\SourceLeaf{363}\ArithmeticHead{32}{L}{STÄCKEL, GAUSS ALS GEOMETER.}\setcounter{footnote}{0}
\noindent Universität Marburg berufen worden war, 1818 mit seinem Kollegen \textsc{Gerling}
darüber gesprochen. »Ich erzählte ihm darauf, wie Sie vor einigen Jahren
[1816] öffentlich geäussert hätten, dass man seit \textsc{Euklids} Zeiten im Grunde
hiermit nicht weiter gekommen sei; ja dass Sie gegen mich mehrmals geäussert
hätten, wie Sie durch vielfältige Beschäftigung mit diesem Gegenstand
auch nicht zum Beweis von der Absurdität einer solchen Annahme [einer
nichteuklidischen Geometrie] gekommen seien« (Brief von \textsc{Gerling} an \textsc{Gauss}
vom 25. Januar 1819, W. VIII, S. 180). \textsc{Schweikart} bat darauf \textsc{Gerling}, er
möge eine kurze Aufzeichnung über seine »Astralische Grössenlehre« (W. VIII,
S. 180) an \textsc{Gauss} weitergeben und diesen ersuchen, ihn gelegentlich sein Urteil
wissen zu lassen. In seiner Antwort erklärt \textsc{Gauss}, es sei ihm fast alles
aus der Seele geschrieben (Brief vom 16. März 1819, W. VIII, S. 181). Er
fand hier die Auffassung wieder, die er in dem Brief an \textsc{Olbers} vom 28. April
1817 ausgesprochen hatte, und die er in dem Brief an \textsc{Bessel} vom 9. April
1830 (W. VIII, S. 201) noch stärker betont hat, dass der Raum eine ausserhalb
von uns vorhandene Wirklichkeit sei, der wir ihre Gesetze nicht \textls[100]{vollständig}
vorschreiben können, deren Eigenschaften vielmehr nur auf Grund
der Erfahrung vollständig festzustellen sind\footnote{Vgl. auch die gegen \textsc{Kant} gerichteten Bemerkungen W. II, S. 177 und W. VIII, S. 224.}.
\par
Das von \textsc{Schweikart} gewählte Beiwort »astralisch« sollte ausdrücken, dass
erst bei Abmessungen der Grössen, wie sie in der Sternenwelt vorkommen,
Abweichungen von der \textsc{Euklid}ischen Geometrie beobachtet werden könnten.
Es scheint \textsc{Gauss} gefallen zu haben, denn er hat es in späteren Aufzeichnungen
angewendet (W. VIII, S. 232).
\par
Ein Neffe von \textsc{Schweikart}, ebenfalls ein Rechtsgelehrter, \textsc{Taurinus} (1794
bis 1874) hatte sich als junger Mann, angeregt durch die Schrift seines Onkels,
mit der Parallelentheorie beschäftigt und im Oktober 1824 einen Beweisversuch
an \textsc{Gauss} gesandt\footnote{Für die folgende Darstellung vgl. P. Th. S. 246—252 und den Aufsatz von P. \textsc{Stäckel}: \textit{F. A. Taurinus}, Abhandlungen zur Geschichte der Mathematik, Heft 9, Leipzig 1899, S. 397.}; dass dieser sich mit den Grundlagen der Geometrie
beschäftige, wusste er seit 1821 durch seinen Onkel. \textsc{Gauss}, der in
\textsc{Taurinus} »einen denkenden mathematischen Kopf« erkannt hatte, antwortete
in einem längeren Schreiben vom 8. November 1824; er hat darin seine Ansichten
über das Parallelenaxiom ausführlich dargelegt, aber zugleich dem
\clearpage\SourceLeaf{364}\ArithmeticHead{33}{R}{DIE GRUNDLAGEN DER GEOMETRIE.}\setcounter{footnote}{0}
\noindent Empfänger des Briefes zur Pflicht gemacht, von dieser »Privat-Mitteilung auf
keine Weise einen öffentlichen oder zur Öffentlichkeit führen könnenden Gebrauch
zu machen« (W. VIII, S. 186—188).
\par
Es muss hier genügen, aus dem Briefe die Hauptstellen anzuführen. »Die
Annahme, dass die Summe der drei Winkel [des Dreiecks] kleiner sei als
\(180^0\), führt auf eine eigene, von der unsrigen (euklidischen) ganz verschiedene
Geometrie, die in sich selbst durchaus konsequent ist und die ich für mich
selbst ganz befriedigend ausgebildet habe, so dass ich jede Aufgabe in derselben
auflösen kann mit Ausnahme der Bestimmung einer Konstante, die
sich a priori nicht ausmitteln lässt. Je grösser man diese Konstante annimmt,
desto mehr nähert man sich der euklidischen Geometrie und ein unendlich
grosser Wert macht beide zusammenfallen.\ .\ .\ .\ Wäre die nichteuklidische
Geometrie die wahre, und jene Konstante in einigem Verhältnisse zu solchen
Grössen, die im Bereich unserer Messungen auf der Erde oder am Himmel
liegen, so liesse sie sich a posteriori ausmitteln«.
\par
Die freundliche Antwort, die der erste Mathematiker der Zeit ihm zukommen
liess, hat \textsc{Taurinus} gewiss angespornt, seine Untersuchungen mit erhöhtem
Eifer fortzusetzen. In seiner 1825 veröffentlichten \textit{Theorie der Parallellinien}
ist er zwar von der unbedingten Gültigkeit des Parallelenaxioms überzeugt,
aber er beginnt die Folgen zu entwickeln, die sich aus dessen Verwerfung
ergeben, und gelangt so seinerseits zu jener Konstanten, die einer nichteuklidischen
Geometrie eigen sein müsste; in der gleichzeitigen Möglichkeit unendlich
vieler solcher Geometrien, die jede für sich genommen widerspruchslos sind,
sieht er jedoch einen ausreichenden Grund, sie alle abzuweisen.
\par
\textsc{Gauss}, dem die Schrift zugesandt wurde, hat sich eben so wenig dazu
geäussert wie zu einer zweiten, den 1826 veröffentlichten \textit{Geometriae prima
elementa}\footnote{Dies geht aus dem Briefe von \textsc{Taurinus} an \textsc{Gauss} vom 29. Dezember 1829 hervor (Brief im \textsc{Gauss}-Archiv). Vermutlich hatte \textsc{Gauss} daran Anstoss genommen, dass er von \textsc{Taurinus} in der Vorrede zur \textit{Theorie der Parallellinien} (S. XIII) und in der Vorrede der \textit{Elementa} (S. V—VI) erwähnt worden war.}. Hier ist \textsc{Taurinus} auf die »neue Geometrie« genauer eingegangen
und hat die Formeln der zugehörigen Trigonometrie sozusagen mit einem
Schlage gewonnen, indem er in den entsprechenden Formeln der sphärischen
Trigonometrie den Halbmesser der Kugel imaginär setzte. Aber noch mehr,
er hat diese Formeln sogleich zur Lösung einer Reihe von Aufgaben angewandt
\par\vfill\noindent{\footnotesize X\textsubscript{2} Abh. 4.\hfill 5\hspace*{2em}}
\clearpage\SourceLeaf{365}\ArithmeticHead{34}{L}{STÄCKEL, GAUSS ALS GEOMETER.}\setcounter{footnote}{0}
\noindent und zum Beispiel den Umfang und den Inhalt des Kreises, die Oberfläche
und das Volumen der Kugel richtig berechnet.
\par
Die Gedanken von \textsc{Taurinus} sind unbeachtet geblieben. »Der Erfolg
bewies mir«, schreibt er am 29. Dezember 1829 an \textsc{Gauss}, »dass Ihre Autorität
dazu gehört, ihnen Anerkennung zu verschaffen, und dieser erste schriftstellerische
Versuch ist, anstatt, wie ich gehofft hatte, mich zu empfehlen, für
mich eine reiche Quelle von Unzufriedenheit geworden« (Brief im \textsc{Gauss}-Archiv).
\par
Im fünften Abschnitt dieses Aufsatzes wird ausführlich über die Untersuchungen
berichtet werden, die \textsc{Gauss} in der Zeit von 1816 bis 1827 über
die allgemeine Lehre von den krummen Flächen angestellt hat. Erst dort
soll auf die Zusammenhänge mit den Grundlagen der Geometrie eingegangen
und im Besonderen die Frage erörtert werden, ob \textsc{Gauss} die Beziehung zwischen
der absoluten Geometrie und der Geometrie auf den Flächen konstanten
Krümmungsmasses gekannt hat. Auch die Ansichten von \textsc{Gauss} über mehrdimensionale
Mannigfaltigkeiten werden dann zur Sprache kommen.
\par
Bald nach der Vollendung der \textit{Disquisitiones generales circa superficies curvas}
(Oktober 1827), die, wie \textsc{Gauss} am 11. Dezember 1825 an \textsc{Hansen} schrieb,
»tief in vieles Andere, ich möchte sogar sagen, in die Metaphysik der Raumlehre
eingreifen« (Brief im \textsc{Gauss}-Archiv), hat sich \textsc{Gauss} erneut den Grundlagen
der Geometrie zugewandt. Am 27. Januar 1829 berichtet er an \textsc{Bessel}:
»Auch über ein anderes Thema, das bei mir fast schon 40 Jahr alt ist, habe
ich zuweilen in einzelnen freien Stunden wieder nachgedacht, ich meine die
ersten Gründe der Geometrie.\ .\ .\ .\ Inzwischen werde ich wohl noch lange
nicht dazu kommen, meine \textls[100]{sehr ausgedehnten} Untersuchungen darüber
zur öffentlichen Bekanntmachung auszuarbeiten, und vielleicht wird dies auch
bei meinen Lebzeiten nie geschehen, da ich das Geschrei der Böoter scheue,
wenn ich meine Ansicht \textls[100]{ganz} aussprechen wollte« (W. VIII, S. 200).
\par
Von den Untersuchungen, auf die \textsc{Gauss} hindeutet, ist uns nur eine kurze
Notiz vom November 1828 erhalten, in der unabhängig vom elften Axiom
bewiesen wird, dass die Winkelsumme des Dreiecks nicht grösser sein kann,
als zwei Rechte (W. VIII, S. 190)\footnote{Dasselbe Verfahren hatte schon \textsc{Legendre} in der zweiten Auflage der \textit{Éléments de géométrie} (1798) angewandt.}. Aber im April 1831 hat er begonnen,
einiges von seinen Meditationen aufzuschreiben. »Ich wünschte doch, dass es
\clearpage\SourceLeaf{366}\ArithmeticHead{35}{R}{DIE GRUNDLAGEN DER GEOMETRIE.}\setcounter{footnote}{0}
\noindent nicht mit mir unterginge« (Brief an \textsc{Schumacher} vom 17. Mai 1831, W. VIII,
S. 213).
\par
Als diese Niederschriften darf man drei Zettel ansprechen, die aus dem
Nachlass W. VIII, S. 202—209 abgedruckt sind. In der Notiz [3], die wohl
die früheste ist, und von der [1] und [2] nur genauere Ausführungen sind\footnote{Nach einer brieflichen Mitteilung von H. S. \textsc{Carslaw} (Sydney) ist in Nr. 4 der Notiz [3] der Fall übersehen, dass die Geraden \(cb\) und {\scriptsize 1} einander nicht schneiden; diese Lücke ist in der Notiz [1], Nr. 4, Fall II ausgefüllt.},
werden die grundlegenden Eigenschaften der parallelen oder, nach \textsc{Johann
Bolyai}, asymptotischen Geraden hergeleitet, und in der letzten Nummer gelangt
\textsc{Gauss} zu dem Parazykel, der Kurve, in die der Kreis übergeht, wenn
der Halbmesser unendlich wird. Er nennt sie Trope, also Wendekreis (cercle
tropique), ein deutliches Zeichen, dass er den Parazykel als den Übergang
von den eigentlichen Kreisen zu den Hyperzykeln aufgefasst hat. Der Gang
der Entwicklung hat grosse Ähnlichkeit mit dem von \textsc{Johann Bolyai} in der
\textit{Scientia spatii}.
\par
Ein Ersatz für weitere Aufzeichnungen, freilich ein spärlicher, ist der
Brief an \textsc{Schumacher} vom 12. Juli 1831, in dem \textsc{Gauss} die Folgen bespricht,
die das Aufhören der Ähnlichkeit in der nichteuklidischen Geometrie nach
sich zieht, und die dort geltende Formel für den Umfang des Kreises angibt
(W. VIII, S. 215). Es ist leider nur wenig, was wir von jenen sehr ausgedehnten
Untersuchungen wissen, und auch aus den folgenden Jahren wird nur
wenig hinzukommen.
\clearpage\endgroup
